\section*{Capítulo \ref{ch:b2:coordination-complexes} --- Metais de transição e complexos de coordenação}

\begin{solution}{exo:b2:coordination-complexes:1}
\ce{Ti^3+} $d^1$, \ce{V^3+} $d^2$, \ce{Cr^3+} $d^3$, \ce{Mn^2+} $d^5$, \ce{Fe^3+} $d^5$, \ce{Co^2+}
$d^7$, \ce{Ni^2+} $d^8$, \ce{Zn^2+} $d^{10}$.
\end{solution}

\begin{solution}{exo:b2:coordination-complexes:2}
Água 1, en 2, etanodioato 2, edta 6, tiocianato 1; o tiocianato é ambidentado
(por S ou por N).
\end{solution}

\begin{solution}{exo:b2:coordination-complexes:3}
Sulfato de tetra-amin\-cobre(II); hexaciani\-doferrato(III) de potássio;
cloreto de tetra-amin\-diclorido\-cobalto(III); tetracarbonil\-níquel(0).
\end{solution}

\begin{solution}{exo:b2:coordination-complexes:4}
Linear; tetraédrica; quadrado plana ($d^8$, terceira série); octaédrica.
\end{solution}

\begin{solution}{exo:b2:coordination-complexes:5}
Dois: cis (os dois cloretos adjacentes) e trans (opostos). Em um tetraedro, só
haveria um.
\end{solution}

\begin{solution}{exo:b2:coordination-complexes:6}
\ce{Cr(CO)6} 18; \ce{Fe(CO)5} 18; \ce{Ni(CO)4} 18; \ce{V(CO)6}: $5 + 12 = 17$;
\ce{Mn2(CO)10} 18 por manganês. \ce{V(CO)6} viola a regra: ele é \omterm{def:b2:diatomic-mos:magnetism}{paramagnético} e
é facilmente reduzido a \ce{[V(CO)6]-}, que tem 18.
\end{solution}

\begin{solution}{exo:b2:coordination-complexes:7}
Ferroceno: Fe(II) $d^6$ + 12 = 18. Cobaltoceno: Co(II) $d^7$ + 12 = 19. O
cobaltoceno perde facilmente seu elétron a mais, dando o íon cobaltocênio
\ce{[Co(C5H5)2]+}, com 18.
\end{solution}

\begin{solution}{exo:b2:coordination-complexes:8}
Um edta substitui seis moléculas de água: \ce{[Ni(H2O)6]^2+ + edta^4- -> [Ni(edta)]^2- +
6H2O}, sete partículas a partir de duas. A \omterm{def:b2:reaction-free-energy:reaction-entropy}{entropia de reação} é fortemente
positiva, e $K$, muito grande. Seis ligantes separados substituem seis
moléculas de água sem ganho no número de partículas.
\end{solution}

\begin{solution}{exo:b2:coordination-complexes:9}
Nitrito-$\kappa N$ (\ce{Co-NO2}, o isômero nitro, amarelo) e nitrito-$\kappa O$
(\ce{Co-O-N=O}, o isômero nitrito, vermelho).
\end{solution}

\begin{solution}{exo:b2:coordination-complexes:10}
Três: o isômero trans (aquiral: planos de simetria) e os dois enantiômeros do
isômero cis (os dois quelatos e os dois cloretos se enrolam em um sentido ou
no outro).
\end{solution}

\begin{solution}{exo:b2:coordination-complexes:11}
Hexágono plano: B nas posições 1,2 (orto), 1,3 (meta), 1,4 (para): três
isômeros. Prisma trigonal: B ao longo de uma aresta de um triângulo, ao longo
de uma aresta vertical, ou em diagonal em uma face retangular: três. O
octaedro dá dois. Não encontrar mais de dois, apesar de muitas tentativas e
de vários compostos, exclui as outras duas geometrias --- se nenhum terceiro
isômero tiver passado despercebido.
\end{solution}

\begin{solution}{exo:b2:coordination-complexes:12}
\ce{[RhCl(PPh3)3]}: Rh(I) $d^8$ + 4 × 2 = 16. \ce{[IrH(CO)(PPh3)3]}: Ir(I) $d^8$ + 5 × 2 =
18. O complexo de ródio, com 16 elétrons, pode receber mais dois (adição
oxidativa, \cref{ch:b2:catalytic-cycles}).
\end{solution}

\begin{solution}{pb:b2:coordination-complexes:1}
\textbf{1.} Os cloretos fora da \omterm{def:b2:coordination-complexes:coordination-sphere}{esfera de coordenação}, livres em solução.
\textbf{2.} A: 0; B: 1; C e D: 2.
\textbf{3.} A: $[\ce{Co(NH3)6}]^{3+}$ + 3 \ce{Cl-}: 4 íons; B: 3; C, D: 2. Eles concordam.
\textbf{4.} +3.
\textbf{5.} $d^6$.
\textbf{6.} Seis em cada um: 6 \ce{NH3}; 5 \ce{NH3} + 1 Cl; 4 \ce{NH3} + 2 Cl.
\textbf{7.} \ce{[Co(NH3)6]Cl3}; \ce{[CoCl(NH3)5]Cl2}; \ce{[CoCl2(NH3)4]Cl} (duas vezes).
\textbf{8.} Cloreto de hexa-amincobalto(III); cloreto de penta-amincloridocobalto(III).
\textbf{9.} Tetra-amindicloridocobalto(III).
\textbf{10.} Isômeros geométricos (cis--trans).
\textbf{11.} \ce{[CoCl3(NH3)3]}, um complexo neutro: nenhum íon, nenhum cloreto
precipitado de imediato.
\textbf{12.} Dois: fac e mer.
\textbf{13.} Cis: os dois Cl a $90^\circ$; trans: a $180^\circ$.
\textbf{14.} Três (posições orto, meta, para no hexágono).
\textbf{15.} Três.
\textbf{16.} Isso sustenta o octaedro, o único dos três que prevê dois.
\textbf{17.} Não: a ausência de um terceiro isômero é uma evidência negativa. Um
terceiro isômero teria excluído o octaedro.
\textbf{18.} O que forma o quelato: um etanodioato bidentado só pode ocupar duas
posições cis.
\textbf{19.} Três quelatos en em torno do cobalto, cada um ocupando duas posições cis.
\textbf{20.} Nenhum dos dois: nem plano de simetria nem centro de inversão.
\textbf{21.} Dois enantiômeros, $\Delta$ e $\Lambda$.
\textbf{22.} Que os seis ligantes estão dispostos em três dimensões, como em um
octaedro; um arranjo plano daria um íon aquiral.
\textbf{23.} Cis: quiral (dois enantiômeros). Trans: aquiral.
\textbf{24.} Ela mostrou que a atividade óptica pertence à geometria molecular em
geral, e não ao carbono, e confirmou o octaedro.
\textbf{25.} O octaedro prevê $\boldsymbol{2}$ isômeros de \ce{[CoCl2(NH3)4]+}; o
hexágono plano e o prisma trigonal preveem $\boldsymbol{3}$.
\end{solution}
